% DevOrbit - Sprint 4 acceptance test report: GitHub integration exit checks (M1-M4).
% Compile:  pdflatex sprint4-exit-checks.tex   (run it twice so the table of contents is filled in)
%
\newcommand{\ResultCheckFive}{\PASS}
\newcommand{\ResultBonus}{\PASS}

\documentclass[11pt,a4paper]{article}
\usepackage[margin=2.2cm]{geometry}
\usepackage[T1]{fontenc}
\IfFileExists{lmodern.sty}{\usepackage{lmodern}}{}
\usepackage{graphicx}
\usepackage{float}
\usepackage{booktabs}
\usepackage{tabularx}
\usepackage{array}
\usepackage[table]{xcolor}
\usepackage{enumitem}
\usepackage{caption}
\usepackage{fancyhdr}
\usepackage[hidelinks]{hyperref}
\usepackage{xurl}

\definecolor{passgreen}{HTML}{1E8E4E}
\definecolor{failred}{HTML}{C62828}
\definecolor{pendorange}{HTML}{E08A00}
\definecolor{headgray}{HTML}{EEF0F5}
\definecolor{accent}{HTML}{6343D7}

\newcommand{\PASS}{\textbf{\textcolor{passgreen}{PASS}}}
\newcommand{\FAIL}{\textbf{\textcolor{failred}{FAIL}}}
\newcommand{\PENDING}{\textbf{\textcolor{pendorange}{PENDING}}}

% \shot{file in images/}{width as fraction of the line}{caption}
\newcommand{\shot}[3]{%
  \begin{figure}[H]
    \centering
    \IfFileExists{images/#1}{%
      \includegraphics[width=#2\linewidth,height=0.40\textheight,keepaspectratio]{images/#1}%
    }{%
      \fcolorbox{gray}{gray!8}{\parbox[c][4.2cm][c]{0.82\linewidth}{\centering\small
        \textbf{Screenshot to add}\\[4pt] save it as \texttt{images/#1}}}%
    }
    \caption{#3}
  \end{figure}}

% One check: objective / requirement / steps / expected / actual / result
\newenvironment{checkinfo}{\par\smallskip\noindent\tabularx{\linewidth}{>{\bfseries}p{2.6cm} X}\toprule}{\bottomrule\endtabularx\par\medskip}

\setlist{itemsep=2pt,topsep=4pt}
\captionsetup{font=small,labelfont=bf}
\pagestyle{fancy}
\fancyhf{}
\lhead{\small DevOrbit -- Sprint 4 acceptance tests}
\rhead{\small GitHub integration (M1--M4)}
\cfoot{\small \thepage}

\hypersetup{pdftitle={DevOrbit Sprint 4 GitHub exit checks}, pdfauthor={Urvish Babariya}}

\begin{document}

\begin{titlepage}
  \centering
  \vspace*{2.5cm}
  {\Huge\bfseries\color{accent} DevOrbit\par}
  \vspace{0.6cm}
  {\LARGE Sprint 4 Acceptance Test Report\par}
  \vspace{0.3cm}
  {\Large GitHub Integration Exit Checks (M1--M4)\par}
  \vspace{1.5cm}
  \begin{tabular}{>{\bfseries}l l}
    User stories   & US-31, US-32, US-33, US-34 \\
    Requirements   & FR-11, FR-12, FR-35, NFR-4, NFR-11a, NFR-11b \\
    Test type      & Manual acceptance (black-box, end-to-end, real GitHub) \\
    Tester         & Urvish Babariya \\
    Date           & 3 November 2026 \\
    Code version   & \texttt{main} at \texttt{b234f22} (PR \#29) \\
    Repository     & \url{https://github.com/Heer64/DevOrbit} \\
  \end{tabular}
  \vfill
  {\small Course team project -- Software Engineering\par}
\end{titlepage}

\tableofcontents
\newpage

%=====================================================================
\section{Purpose and scope}
The Sprint 4 plan ends with exit checks for the GitHub integration that the automated tests cannot cover,
because they need a real GitHub account, a real OAuth App and a real repository sending webhooks.
This report records those checks, run by hand against the code on \texttt{main} after all eight
Sprint 4 parts were merged.

\begin{itemize}
  \item \textbf{M1, US-31} Authorize GitHub (OAuth, state check, token encrypted, never sent to the browser).
  \item \textbf{M2, US-32} Connect one repository; the connection shows as Active.
  \item \textbf{M3, US-33} Owner/Admin view status, disconnect, re-authorize; old feed entries stay.
  \item \textbf{M4, US-34} Signed webhooks for commits, issues and pull requests appear in the Activity feed.
\end{itemize}

The other Sprint 4 parts (M5 files, M6 docs, M7 dashboard, M8 completion and health) are covered by the
automated backend, frontend and Playwright end-to-end tests and are not repeated here.

%=====================================================================
\section{Summary of results}
\begin{center}
\rowcolors{2}{white}{headgray}
\begin{tabularx}{\linewidth}{l X l l}
  \toprule
  \textbf{\#} & \textbf{Exit check} & \textbf{Stories} & \textbf{Result} \\
  \midrule
  1  & Demo: connect GitHub and a repository (Active, webhook active, ping received) & US-31, US-32, US-34 & \PASS \\
  1b & A repository the account does not administer: connected, but no webhook and a clear message & US-32, US-34 & \PASS \\
  2  & A real commit and a real issue appear in the Activity feed with a GitHub link & US-34 & \PASS \\
  3  & A webhook request with a wrong or missing signature is rejected (401) and nothing is stored & US-34, NFR-4 & \PASS \\
  4  & The app keeps working while GitHub is unreachable; GitHub actions fail safely & NFR-11b & \PASS \\
  5  & Disconnect: webhook removed, new events stop, old entries stay & US-33 & \ResultCheckFive \\
  B  & Bonus: re-authorize after disconnect, connection Active again & US-33 & \ResultBonus \\
  \bottomrule
\end{tabularx}
\end{center}

%=====================================================================
\section{Test environment}
\begin{center}
\begin{tabularx}{\linewidth}{>{\bfseries}p{3.6cm} X}
  \toprule
  Machine        & Windows laptop, Google Chrome, VS Code terminals (PowerShell) \\
  Backend        & \texttt{backend/}, \texttt{npm run dev} on \path{http://localhost:5000} \\
  Frontend       & \texttt{frontend/}, \texttt{npm run dev} on \path{http://localhost:5173} \\
  Database       & MongoDB Community Server 9.0.4 installed locally (\path{mongodb://127.0.0.1:27017/devorbit}) \\
  GitHub OAuth App & ``DevOrbit Local Test'', callback \path{http://localhost:5000/api/github/callback} \\
  Webhook relay  & smee.io channel + \texttt{smee-client}, forwarding to \texttt{/api/github/webhook} (GitHub cannot reach \texttt{localhost}) \\
  GitHub account & \texttt{@Urvish2007} \\
  Test repository & \path{Urvish2007/devorbit-webhook-test} (public, owned by the tester, so the account is its admin) \\
  Secrets        & \texttt{GITHUB\_CLIENT\_ID}, \texttt{GITHUB\_CLIENT\_SECRET}, \texttt{GITHUB\_TOKEN\_KEY} (64 hex), \texttt{GITHUB\_WEBHOOK\_SECRET}, \texttt{GITHUB\_WEBHOOK\_URL} set in \texttt{backend/.env} only; never committed and not shown in this report \\
  \bottomrule
\end{tabularx}
\end{center}

\subsection{Setup}
\begin{enumerate}
  \item Registered the OAuth App on GitHub with the homepage and callback URLs above (Figure~\ref{fig:oauth}).
  \item Created a smee.io channel and started \texttt{smee-client} to forward deliveries to the local API.
  \item Filled in the GitHub settings in \texttt{backend/.env}; the token key and webhook secret were generated with
        \path{crypto.randomBytes(32)}.
  \item Before the settings were filled in, the server started normally and listed only the \emph{names} of the
        missing settings, never any values (Figure~\ref{fig:before}). This matches US-31: secrets live only in server
        environment variables and are never logged.
\end{enumerate}

\shot{setup-oauth-app.png}{0.85}{GitHub OAuth App registration: homepage \texttt{http://localhost:5173}, redirect URI \texttt{http://localhost:5000/api/github/callback}, wildcard matching and device flow off.\label{fig:oauth}}
\shot{setup-backend-before-config.png}{0.6}{Backend before the GitHub settings were added: GitHub connect is turned off and only the setting names are printed. The API still starts and connects to MongoDB.\label{fig:before}}

%=====================================================================
\section{Exit checks}

%---------------------------------------------------------------------
\subsection{Check 1: Demo, connect GitHub and a repository}
\begin{checkinfo}
  Objective   & The Owner authorizes GitHub, picks one repository, the connection shows as Active and the webhook is created. \\
  Requirement & US-31 (authorize with a state-checked callback), US-32 (one repository, shown as Active), US-34 (signed webhook configured). \\
  Steps       & Signed up, created a project, clicked \emph{Connect GitHub}, approved on GitHub, clicked \emph{Choose repository},
                selected \path{Urvish2007/devorbit-webhook-test}, clicked \emph{Connect repository}. \\
  Expected    & Panel shows \emph{Active} and \emph{Webhook active}; GitHub sends a \texttt{ping} that DevOrbit records. \\
  Actual      & Authorized as \texttt{@Urvish2007}; repository connected; status \emph{Active}; badge \emph{Webhook active}.
                smee.io received the \texttt{ping} from \path{Urvish2007/devorbit-webhook-test}, and \emph{Recent deliveries}
                lists it as \emph{Ignored} (``GitHub ping: the webhook reached DevOrbit''), which is the intended handling of a ping. \\
  Result      & \PASS \\
\end{checkinfo}
\shot{check1-webhook-active.png}{0.95}{GitHub panel after connecting \texttt{Urvish2007/devorbit-webhook-test}: connection status Active, authorized as \texttt{@Urvish2007}, webhook active.}
\shot{check1-smee-ping.png}{0.6}{smee.io received GitHub's \texttt{ping} event from \texttt{Urvish2007/devorbit-webhook-test} when the webhook was created.}

\subsubsection*{Check 1b: repository the account does not administer}
\begin{checkinfo}
  Objective   & GitHub only lets repository admins add webhooks. Connecting someone else's repository must not fail silently. \\
  Steps       & In the first attempt, the repository \path{NiranjanD0/WarpCopy-Parallel-File-Transfer-over-SSH} (owned by another account) was chosen. \\
  Expected    & The repository connects (US-32), but no webhook is created, and the panel explains why. \\
  Actual      & Status \emph{Active}; badge \emph{Webhook not set up} with the message ``the authorized GitHub account is not an admin of
                this repository, so it cannot add a webhook''. No deliveries. Switching with \emph{Change repository} to the tester's own
                repository then worked (Check 1). \\
  Result      & \PASS \\
\end{checkinfo}
\shot{check1b-not-admin.png}{0.95}{Non-admin repository: connected and Active, but the webhook is not set up and the reason is shown.}

%---------------------------------------------------------------------
\subsection{Check 2: Real GitHub events appear in the Activity feed}
\begin{checkinfo}
  Objective   & A commit pushed, or an issue opened, in the connected repository appears in the project's Activity feed with
                event type, title, author and a link to GitHub, within 60 seconds. \\
  Requirement & US-34, FR-12. \\
  Steps       & On GitHub: edited \texttt{README.md} and committed to \texttt{main}; opened issue \#1 ``Test issue from exit check''.
                In DevOrbit: opened the \emph{Activity} tab and \emph{Recent deliveries}. \\
  Expected    & Both events delivered (HTTP 202), processed, and shown in the feed with an \emph{Open on GitHub} link. \\
  Actual      & smee forwarded three deliveries: \texttt{ping} (200) then \texttt{push} and \texttt{issues} (202 each).
                \emph{Recent deliveries}: \texttt{push} and \texttt{issues~·~opened} both ``Added to feed, 1 new entry''.
                Activity feed: ``Urvish2007 pushed commit `Update README with webhook test information' to main'' and
                ``Urvish2007 opened issue \#1 `Test issue from exit check''', each with \emph{Open on GitHub}.
                Deliveries are answered at once and processed straight away, well inside the 60-second limit. \\
  Result      & \PASS \\
\end{checkinfo}
\shot{check2-smee-terminal.png}{0.55}{smee-client forwarding GitHub's deliveries to the local API: \texttt{ping} answered 200, \texttt{push} and \texttt{issues} accepted with 202.}
\shot{check2-recent-deliveries.png}{0.95}{Recent deliveries in the GitHub panel: \texttt{issues~·~opened} and \texttt{push} added to the feed, \texttt{ping} ignored.}
\shot{check2-activity-feed.png}{0.8}{Activity feed entries created from GitHub, each with type, title, GitHub author and an \emph{Open on GitHub} link.}

%---------------------------------------------------------------------
\subsection{Check 3: Requests with a bad or missing signature are rejected}
\begin{checkinfo}
  Objective   & Only deliveries signed with the shared secret are accepted (\texttt{X-Hub-Signature-256}, HMAC-SHA256 of the raw body). \\
  Requirement & US-34 (verify the signature before processing), NFR-4. \\
  Steps       & Sent a fake \texttt{push} delivery with \texttt{curl.exe}: (a) with signature \texttt{sha256=000\ldots0}
                (delivery \texttt{manual-test-1}); (b) with no signature header (delivery \texttt{manual-test-2}). \\
  Expected    & HTTP 401 for both; nothing stored in the feed or in Recent deliveries. \\
  Actual      & Both requests returned \textbf{401}. The Activity feed and Recent deliveries did not change; neither delivery ID was recorded. \\
  Result      & \PASS \\
\end{checkinfo}
\shot{check3-no-signature.png}{0.6}{Fake deliveries rejected: wrong signature \texttt{sha256=000\ldots0} returns 401, and a request without a signature returns 401.}

%---------------------------------------------------------------------
\subsection{Check 4: The app works while GitHub is unreachable}
\begin{checkinfo}
  Objective   & An outage at GitHub must not affect the rest of DevOrbit, and GitHub actions must fail without changing data. \\
  Requirement & NFR-11b; Sprint 4 exit check ``app works with GitHub down''. \\
  Steps       & Blocked GitHub on the test PC by adding \path{127.0.0.1 github.com} and \path{127.0.0.1 api.github.com} to the
                Windows \texttt{hosts} file. Confirmed with \path{curl.exe https://api.github.com} (result \texttt{000}, no connection).
                Then used the project: created a task, moved it on the board, sent a chat message, created a docs page, opened the Activity feed.
                Clicked \emph{Check connection}. Finally removed the two lines and clicked \emph{Check connection} again. \\
  Expected    & All project features work. \emph{Check connection} shows an error and the status stays Active. After GitHub is reachable again,
                the check succeeds. \\
  Actual      & Task ``Task while GitHub is down'' created and moved To do $\rightarrow$ In progress; chat message ``Chat works without GitHub''
                posted in \texttt{\#general}; docs page ``Offline note'' created; all three logged in the Activity feed.
                \emph{Check connection} showed ``Could not reach GitHub. Please try again.'' and the status stayed \emph{Active}, with the webhook still active.
                After unblocking: ``Connection checked: the authorization works and \path{Urvish2007/devorbit-webhook-test} can be read.'' \\
  Result      & \PASS \\
\end{checkinfo}
\shot{check4-github-blocked.png}{0.55}{GitHub blocked: \texttt{curl.exe https://api.github.com} returns \texttt{000} (cannot connect).}
\shot{check4-board.png}{0.65}{While GitHub is down: a new task created and moved to In Progress on the Kanban board.}
\shot{check4-chat.png}{0.6}{While GitHub is down: chat message posted in \texttt{\#general}.}
\shot{check4-docs.png}{0.6}{While GitHub is down: documentation page ``Offline note'' created.}
\shot{check4-activity.png}{0.65}{While GitHub is down: the Activity feed records the new task, the move and the chat message.}
\shot{check4-check-connection-error.png}{0.6}{\emph{Check connection} while GitHub is down: a clear error, and the connection stays Active with the webhook still active (nothing changed).}
\shot{check4-check-connection-ok.png}{0.6}{After GitHub was unblocked: \emph{Check connection} succeeds.}

%---------------------------------------------------------------------
\subsection{Check 5: Disconnect, new events stop and old entries stay}
\begin{checkinfo}
  Objective   & After disconnect, no new GitHub events reach the feed, the webhook is removed from GitHub, and earlier feed entries remain. \\
  Requirement & US-33, FR-35. \\
  Steps       & Clicked \emph{Disconnect} and confirmed. While disconnected, on GitHub: committed the README line
                ``After disconnect -- should NOT appear'' and opened issue \#2 ``Opened issue after Disconnect!''.
                Then checked the Activity feed. \\
  Expected    & Status \emph{Disconnected}; no feed entry for the commit or the issue made while disconnected; a ``disconnected GitHub''
                entry in the feed; earlier entries are not deleted. \\
  Actual      & Status \emph{Disconnected}: ``GitHub was disconnected from this project. No new GitHub events are added to the Activity Feed.
                Earlier entries stay.'' The panel keeps the last account and repository and offers \emph{Re-authorize GitHub}.
                Both events exist on GitHub (Figures~\ref{fig:readme} and~\ref{fig:issues}), but the feed has \textbf{no entry for either}:
                directly after ``Urvish Babariya disconnected GitHub (repository `Urvish2007/devorbit-webhook-test')'' the next entry is the
                re-authorization (Figure~\ref{fig:after}). The events from while it was disconnected were not added later either.
                Disconnect deletes no feed entries; the screenshot is cropped to the newest entries. The automated test
                \path{githubConnectionManagement.test.js} also checks that the webhook is removed from GitHub on disconnect;
                GitHub's Webhooks settings page was not captured in this run. \\
  Result      & \ResultCheckFive \\
\end{checkinfo}
\shot{check5-disconnected.png}{0.6}{GitHub panel after Disconnect: status Disconnected, last account and repository kept, \emph{Re-authorize GitHub} offered.}
\shot{check5-github-readme.png}{0.95}{On GitHub, \texttt{README.md} contains the line ``After disconnect -- should NOT appear'', committed while DevOrbit was disconnected (later lines were added after re-authorizing).\label{fig:readme}}
\shot{check5-github-issues.png}{0.95}{On GitHub, issue \#2 ``Opened issue after Disconnect!'' was opened while DevOrbit was disconnected.\label{fig:issues}}
\shot{check5-activity-after-disconnect.png}{0.85}{DevOrbit Activity feed, newest first: the commit after re-authorizing, ``reconnected GitHub as @Urvish2007'', and ``disconnected GitHub''. Nothing from the time it was disconnected (no README commit, no issue \#2) was added.\label{fig:after}}

%---------------------------------------------------------------------
\subsection{Bonus: Re-authorize after disconnect}
\begin{checkinfo}
  Objective   & The Owner (or an Admin) can re-authorize a disconnected integration and it returns to Active with its repository. \\
  Requirement & US-33. \\
  Steps       & Clicked \emph{Re-authorize GitHub}, approved on GitHub, then committed ``Update README with reauthorisation note''
                (README line ``we're reauthorising again!''). \\
  Expected    & The connection becomes Active again with the same repository, the re-authorization is logged, and new events reach the feed again. \\
  Actual      & The feed shows ``Urvish Babariya reconnected GitHub as @Urvish2007'', followed about a minute later by
                ``Urvish2007 pushed commit `Update README with reauthorisation note' to main'' with \emph{Open on GitHub} (``just now'').
                So the integration is back on the same repository and its webhook delivers again. \\
  Result      & \ResultBonus \\
\end{checkinfo}
\noindent See Figure~\ref{fig:after}: the ``reconnected GitHub'' entry and the commit that followed it.

%=====================================================================
\section{Observations and defects found}
\begin{center}
\rowcolors{2}{white}{headgray}
\begin{tabularx}{\linewidth}{l l X X}
  \toprule
  \textbf{ID} & \textbf{Severity} & \textbf{Observation} & \textbf{Recommendation} \\
  \midrule
  D1 & Medium &
  When deliveries arrived through smee (which adds an \texttt{X-Forwarded-For} header), the backend logged
  \path{ERR_ERL_UNEXPECTED_X_FORWARDED_FOR} from \texttt{express-rate-limit} (Figure~\ref{fig:d1}). Deliveries were still processed (202).
  The app does not set Express's \texttt{trust proxy}. On Render, which also sits behind a proxy, every client would then share the proxy's
  IP for rate limiting. &
  Set \path{app.set('trust proxy', 1)} in \texttt{createApp()} for the deployed API (one proxy hop), ideally controlled by an environment variable, and re-test. \\
  D2 & Low &
  The Activity item for an issue that was \emph{opened} is titled ``GitHub issue updated'' (its text correctly says ``opened issue \#1''). &
  Title the item by action, e.g.\ ``GitHub issue opened'', or use a neutral ``GitHub issue''. \\
  O1 & Info &
  The OAuth App form had ``Expire user access tokens'' ticked. With it on, GitHub's token expires after 8 hours. The backend stores only the access
  token (no refresh), so the panel would then ask to re-authorize (handled by US-33). &
  For the team's OAuth App, untick it, or plan token refresh as future work. \\
  O2 & Info &
  On the test PC \texttt{ipconfig /flushdns} needs an administrator terminal (Figure~\ref{fig:flush}). Windows applied the \texttt{hosts} change anyway
  (confirmed by \texttt{curl}, Figure in Check 4). &
  Note it in the test procedure; it does not affect the app. \\
  \bottomrule
\end{tabularx}
\end{center}
\shot{obs-rate-limit-warning.png}{0.95}{D1: \texttt{express-rate-limit} validation error logged while smee forwarded deliveries with an \texttt{X-Forwarded-For} header.\label{fig:d1}}
\shot{check4-flushdns-elevation.png}{0.5}{O2: \texttt{ipconfig /flushdns} requires an elevated (administrator) terminal on the test PC.\label{fig:flush}}

%=====================================================================
\section{Related automated tests}
The same rules are also covered by automated tests on \texttt{main}, which run in GitHub Actions on every pull request
(Node 20 and 22, MongoDB 7). They use a fake GitHub, so this manual run is what proves the real integration.

\begin{center}
\rowcolors{2}{white}{headgray}
\begin{tabularx}{\linewidth}{l X X}
  \toprule
  \textbf{Part} & \textbf{Backend (node:test + supertest)} & \textbf{Frontend (Vitest + Testing Library)} \\
  \midrule
  M1 US-31 & \texttt{githubOAuth.test.js}, \path{githubCrypto.test.js} & \path{GithubIntegrationPanel.test.jsx} \\
  M2 US-32 & \texttt{githubRepos.test.js}, \path{githubRepository.test.js} & \path{GithubRepository.test.jsx} \\
  M3 US-33 & \path{githubConnection.test.js}, \path{githubConnectionManagement.test.js} & \path{GithubConnectionManagement.test.jsx} \\
  M4 US-34 & \path{githubEvents.test.js}, \path{githubWebhook.test.js}, \path{githubWebhooks.test.js} & \path{GithubEvents.test.jsx} \\
  \bottomrule
\end{tabularx}
\end{center}

Examples of what they check: a missing or wrong \texttt{X-Hub-Signature-256} returns 401 and stores nothing; a repeated delivery is not
processed twice; after disconnect, new events stop and old entries stay; disconnect and \emph{Check connection} behave safely when GitHub is down
or rate-limited; the token never appears in a response or a log line.

%=====================================================================
\section{Conclusion}
Checks 1--4 passed against real GitHub: authorization, repository connection, webhook creation, delivery of real commit and issue events
into the Activity feed, rejection of unsigned or wrongly signed requests, and correct behaviour while GitHub is unreachable.
Disconnect stopped new events (a commit and an issue made while disconnected never reached the feed), and re-authorizing restored the
connection on the same repository, with new commits arriving again. All Sprint 4 GitHub exit checks passed.
Defect D1 (\texttt{trust proxy}) should be fixed before the API is relied on behind Render's proxy; D2 is a wording issue.

\vspace{1.2cm}
\noindent
\begin{tabularx}{\linewidth}{X X}
  Tested by: \rule{5cm}{0.4pt} & Reviewed by: \rule{5cm}{0.4pt} \\[2pt]
  Urvish Babariya (tester) & \\[14pt]
  Date: 3 November 2026 & Date: \rule{3.5cm}{0.4pt} \\
\end{tabularx}

\end{document}
